\section{Conclusion and Limitations}
\label{sec:conclusion}

\subsection{Conclusion}
This research addressed the complex challenge of integrating hydro-meteorological variables, temporal sequence information, and leakage-aware evaluation frameworks to develop a multi-hazard climate risk prediction and early-warning system. ClimateGuardian AI demonstrates an integrated methodology comprising a 45-feature hydro-climate representation, strict chronological partitioning, leakage-aware preprocessing, and explicit heatwave proxy-feature remediation. By securely connecting these frozen scientific artifacts to a realtime Open-Meteo interface, the framework bridges historical evaluation and operational execution.

The experimental evidence establishes distinct performance profiles across hazards. The Flood Random Forest model provides moderate discriminatory performance ($F_1 = 0.4686$, ROC-AUC $= 0.7184$) governed by a recall-oriented trade-off. The leakage-remediated Heatwave LightGBM model retains strong predictive discrimination ($F_1 = 0.8878$, ROC-AUC $= 0.9969$) without relying on deterministic proxy features. For Compound hazard prediction, the temporal GRU provides the strongest verified offline benchmark ($F_1 = 0.4343$, ROC-AUC $= 0.9598$). However, a Logistic Regression Stacking Ensemble ($F_1 = 0.3419$, ROC-AUC $= 0.9495$) was retained for the operational pipeline due to its critical compatibility with the stateless, point-in-time architecture required for live execution.

Supplementary diagnostics provided crucial insights under the evaluated configurations. Ablation analysis indicated a substantial contribution from temporal variables, increasing Compound $F_1$ from 0.2400 to 0.4051. Platt scaling calibration substantially reduced the Compound Brier score from 0.2013 to 0.0338. While realtime testing successfully validated end-to-end live inference capabilities at selected geographic coordinates, predictive performance relies exclusively on the held-out 2016--2018 dataset.

\subsection{Limitations and Future Work}
While the evaluated framework establishes a robust operational architecture, several scientific limitations dictate directions for future work. First, broader geographic generalization remains to be established, as models were evaluated on a specific regional historical dataset. Second, predicting Compound risk remains difficult; the operational Stacking Ensemble demonstrates lower offline $F_1$ performance than the GRU benchmark. Third, because the current operational pipeline prioritizes point-in-time stateless inference, it does not currently exploit persistent recurrent states. Fourth, predictive performance inherently degrades at extended horizons, as quantified by the Heatwave lead-time results. Finally, the feature space is constrained by the available temporal resolution.

Future work will prioritize broader geographic validation, transitioning to spatial or gridded modeling, and integrating richer hydrological, topographic, and land-surface variables. Enhancing the live architecture to support stateful sequence models could significantly improve realtime Compound-event prediction. Expanding research to larger multi-region datasets, incorporating uncertainty-aware forecasting, and pursuing additional external observational validation will further advance the resilience of multi-hazard predictive systems.
